\section*{Capítulo \ref{ch:g12:reaction-rates} --- Velocidad de reacción}

\begin{solution}{exo:g12:reaction-rates:1}
Temperatura; superficie de contacto; concentración; temperatura.
\end{solution}

\begin{solution}{exo:g12:reaction-rates:2}
\qty{13.9}{min} (azul) y \qty{6.9}{min} (rojo). Al cabo de dos \omterm{def:g12:reaction-rates:half-life}{tiempos
de semirreacción},
$10.0 / 4 = \qty{2.5}{mmol/L}$.
\end{solution}

\begin{solution}{exo:g12:reaction-rates:3}
Pendiente $(5.0 - 1.0)/10 = \qty{0.40}{mmol.L^{-1}.min^{-1}}$: es la
\omterm{def:g12:reaction-rates:rate}{velocidad de aparición} a los \qty{5}{min}.
\end{solution}

\begin{solution}{exo:g12:reaction-rates:4}
Enfriar y diluir hacen la reacción tan lenta que la composición ya no
cambia durante la \omterm{def:g11:titration:titration}{valoración}. La composición medida es la que tenía la
\omterm{def:g6:pure-substances-and-mixtures:mixture}{mezcla} cuando se detuvo la reacción, en el momento de verterla.
\end{solution}

\begin{solution}{exo:g12:reaction-rates:5}
\qty{30}{min} son 3 \omterm{def:g12:reaction-rates:half-life}{tiempos de semirreacción}: queda $1/8$. \qty{1}{h}
son 6 \omterm{def:g12:reaction-rates:half-life}{tiempos de semirreacción}: $1/64$, alrededor del
\qty{1.6}{\percent}.
\end{solution}

\begin{solution}{exo:g12:reaction-rates:6}
$\ln[\ce{A}]$: 2.079, 1.733, 1.386, 1.040, 0.693. Disminuye 0.346 cada
\qty{5}{min}: una recta, primer orden. $k = 0.346/5 =
\qty{0.0693}{min^{-1}}$ y $t_{1/2} = \ln 2 / k = \qty{10.0}{min}$ (como
se ve directamente: 8.00, 4.00, 2.00 cada \qty{10}{min}).
\end{solution}

\begin{solution}{exo:g12:reaction-rates:7}
$k = 0.693 / 25 = \qty{0.0277}{s^{-1}}$.
\end{solution}

\begin{solution}{exo:g12:reaction-rates:8}
$0.100\,e^{-0.20} = \qty{0.0819}{mol/L}$; $0.100\,e^{-0.60} =
\qty{0.0549}{mol/L}$; $0.100\,e^{-2.0} = \qty{0.0135}{mol/L}$.
\end{solution}

\begin{solution}{exo:g12:reaction-rates:9}
A los \qty{13.9}{min}, $[\ce{A}] = \qty{5.0}{mmol/L}$; la pendiente de la
tangente es próxima a $-0.25$, y $-k[\ce{A}] = -0.050 \times 5.0 =
\qty{-0.25}{mmol.L^{-1}.min^{-1}}$: la mitad del valor inicial.
\end{solution}

\begin{solution}{exo:g12:reaction-rates:10}
El yodo es la única especie coloreada: por la ley de Beer--Lambert,
$A = \varepsilon \ell [\ce{I2}]$. La \omterm{def:g12:reaction-rates:rate}{velocidad de aparición} es la
pendiente de la tangente a la curva $[\ce{I2}](t)$, es decir, la
pendiente de la tangente a $A(t)$ dividida por $\varepsilon \ell$.
\end{solution}

\begin{solution}{exo:g12:reaction-rates:11}
La velocidad inicial $k[\ce{A}]_0$ es el doble en el segundo ensayo; los
\omterm{def:g12:reaction-rates:half-life}{tiempos de semirreacción} son iguales, $\ln 2 / k$, independientes de
$[\ce{A}]_0$.
\end{solution}

\begin{solution}{exo:g12:reaction-rates:12}
$t = \ln(1/f)/k$. En \omterm{def:g12:reaction-rates:half-life}{tiempos de semirreacción}: $\ln(1/f)/\ln 2$. Para el
\qty{1}{\percent}, $\ln 100 / \ln 2 = 6.6$; para el \qty{0.1}{\percent},
$\ln 1000 / \ln 2 =
10.0$.
\end{solution}

\begin{solution}{exo:g12:reaction-rates:13}
\qty{90}{\percent} consumido, $f = 0.10$: $t = \ln 10 / k$, es decir,
\qty{46}{min} a \qty{20}{\celsius} y \qty{23}{min} a \qty{30}{\celsius}.
La \omterm{def:g12:reaction-rates:first-order}{constante de velocidad} se duplica en \qty{10}{\celsius}: aquí se
cumple la regla.
\end{solution}

\begin{solution}{exo:g12:reaction-rates:14}
La temperatura baja hasta cerca de \qty{0}{\celsius}, y las
concentraciones se dividen por 10. Con una velocidad proporcional al
producto de dos concentraciones, la \omterm{def:g10:concentration-and-dilution:dilution}{dilución} por sí sola la divide por
$10 \times 10 = 100$.
\end{solution}

\begin{solution}{exo:g12:reaction-rates:15}
$v(t) = -\dfrac{d}{dt}\big([\ce{A}]_0 e^{-kt}\big) = k[\ce{A}]_0 e^{-kt}$.
Velocidad inicial $k[\ce{A}]_0$; la velocidad se reduce a la mitad cuando
$e^{-kt} = 1/2$, es decir, al cabo de un \omterm{def:g12:reaction-rates:half-life}{tiempo de semirreacción}.
\end{solution}

\begin{solution}{pb:g12:reaction-rates:1}
\textbf{1.} Por la ley de Beer--Lambert, ya que el colorante es la única
especie que absorbe a esta longitud de onda; la disolución debe estar lo
bastante diluida ($A$ por debajo de 1, aproximadamente).

\textbf{2.} No absorbe luz visible: no añade nada a $A$.

\textbf{3.} Para que su concentración apenas cambie: así la velocidad
solo depende del colorante.

\textbf{4.} $A = 0$ (el blanco).

\textbf{5.} $A$ baja de 0.800 a 0.402 en \qty{6}{min}: unos
\qty{6}{min}.

\textbf{6.} De 0.402 a los \qty{6}{min} a 0.199 a los \qty{12}{min}: de
nuevo se reduce a la mitad en \qty{6}{min}.

\textbf{7.} $-0.223$, $-0.448$, $-0.691$, $-0.911$, $-1.155$, $-1.366$,
$-1.614$, $-1.945$, $-2.551$, $-3.079$.

\textbf{8.} Los puntos están alineados: la reacción es de primer orden
respecto al colorante.

\textbf{9.} $(-2.551 + 0.223)/20 = \qty{-0.116}{min^{-1}}$.

\textbf{10.} $k = \qty{0.116}{min^{-1}}$.

\textbf{11.} $0.693 / 0.116 = \qty{6.0}{min}$, como se lee en la tabla.

\textbf{12.} $k A_0 = 0.116 \times 0.800 = \qty{0.093}{min^{-1}}$, en
unidades de \omterm{def:g11:absorbance:absorbance}{absorbancia} por minuto.

\textbf{13.} El mismo, \qty{6.0}{min}: el \omterm{def:g12:reaction-rates:half-life}{tiempo de semirreacción} de
primer orden no depende del valor inicial.

\textbf{14.} $0.800\, e^{-0.116 \times 30} = 0.025$.

\textbf{15.} $t = \ln 100 / k = 4.61 / 0.116 = \qty{40}{min}$.

\textbf{16.} $40 / 6.0 = 6.6$ \omterm{def:g12:reaction-rates:half-life}{tiempos de semirreacción}.

\textbf{17.} $t = \ln(0.800/0.010)/k = 4.38 / 0.116 = \qty{38}{min}$.

\textbf{18.} Siguen siendo 6.6 \omterm{def:g12:reaction-rates:half-life}{tiempos de semirreacción}, ahora
$6.6 \times 3.0 = \qty{20}{min}$.

\textbf{19.} Unos 6.6 \omterm{def:g12:reaction-rates:half-life}{tiempos de semirreacción}, es decir, unos
\qty{40}{min} aquí.
\end{solution}
